\section{Lesson 16 --- 26/11/2026}

\subsection{Exact penalty method}
The penalty method of the previous lesson finds an optimal solution only in the limit $\varepsilon_k \to 0$, and small values of $\varepsilon$ cause numerical issues. We now see a variant that stops after a finite number of iterations. The difference from the quadratic penalty method is the definition of the penalty function, while the problem and the considerations remain the same. Consider a convex constrained problem
\[
\begin{cases}
\min f(x)\\
g_i(x) \le 0 & \forall\, i=1,\dots,m.
\end{cases}
\tag{P}
\]
We define the \g{linear penalty function}
\[
\tilde p(x) = \sum_{i=1}^{m} \max\{0, g_i(x)\}
\]
and consider the penalized problem
\begin{equation}\label{eq:exact-Peps}
\begin{cases}
\displaystyle \min f(x) + \frac{1}{\varepsilon}\, \tilde p(x) := \tilde p_\varepsilon(x)\\[1ex]
x \in \R^n,
\end{cases}
\tag{$\tilde P_\varepsilon$}
\end{equation}
which is unconstrained, convex and nonsmooth. As before,
\[
\tilde p_\varepsilon(x)
\begin{cases}
= f(x) & \text{if } x \in X,\\
> f(x) & \text{if } x \notin X.
\end{cases}
\]

Using a linear penalty function we lose differentiability. Luckily for us, there are methods to solve convex nondifferentiable problems, but we will not see them.

For such penalized problem we do not need a sequence $\varepsilon_k \to 0$ to approximate an optimal solution of (P), which avoids numerical issues. In fact, there exists a suitable $\varepsilon$ such that the minimum of \eqref{eq:exact-Peps} coincides with the minimum of (P).

\begin{proposition}
Suppose that there exists an optimal solution $x^*$ of (P) and that $\lambda^*$ is a KKT multipliers vector associated to $x^*$. Then, the sets of optimal solutions of (P) and \eqref{eq:exact-Peps} coincide provided that $\varepsilon \in (0, 1/\norm{\lambda^*}_\infty)$.
\end{proposition}

Since $\lambda^*$ is not known in advance, we cannot choose such $\varepsilon$ directly. We reduce $\varepsilon$ as in the penalty method, but now the threshold is reached after finitely many steps.

\begin{algorithm}[htbp]
\caption{Exact penalty method}
Set $\varepsilon_0 > 0$, $\tau \in (0,1)$, $k = 0$\;
Find an optimal solution $x^k$ of the penalized problem $(\tilde P_{\varepsilon_k})$\;
\While{$x^k \notin X$}{
    $\varepsilon_{k+1} = \tau\, \varepsilon_k$\;
    $k = k + 1$\;
    Find an optimal solution $x^k$ of the penalized problem $(\tilde P_{\varepsilon_k})$\;
}
\end{algorithm}

\begin{theorem}
The exact penalty method stops after a finite number of iterations at an optimal solution of (P).
\end{theorem}

Notice that penalty methods generate a sequence of unfeasible points that approximate an optimal solution of (P).

\subsection{Barrier methods}

Unlike penalty methods, barrier methods generate a sequence of feasible points that converge to an optimal solution of the problem. We need to find a suitable unconstrained problem that approximates our constrained problem on the interior of the feasible set. Consider the problem
\[
\begin{cases}
\min f(x)\\
g_i(x) \le 0 & \forall\, i=1,\dots,m.
\end{cases}
\tag{P}
\]
We need to make some assumptions:
\begin{itemize}
    \item $f$ and $g_i$ are convex and twice continuously differentiable (on $\operatorname{int}(X) \ne \emptyset$);
    \item there exists an optimal solution (e.g.\ $f$ is coercive or $X$ is bounded);
    \item the Slater constraint qualification holds: there exists $\bar x$ such that
    \[
    g_i(\bar x) < 0 \qquad \forall\, i=1,\dots,m.
    \]
\end{itemize}
Hence strong duality holds. Special cases are linear programming and convex quadratic programming.

\subsection{Logarithmic barrier}

On the interior $\operatorname{int}(X)$ of the feasible set $X$, we can approximate the given problem (P) with
\begin{equation}\label{eq:barrier}
\begin{cases}
\displaystyle \min \psi_\varepsilon(x) := f(x) - \varepsilon \sum_{i=1}^{m} \log(-g_i(x))\\[1ex]
x \in \operatorname{int}(X).
\end{cases}
\end{equation}

\begin{definition}[Logarithmic barrier function]
The \g{logarithmic barrier function} is
\[
B(x) := -\sum_{i=1}^{m} \log(-g_i(x)).
\]
\end{definition}
Then
\[
\psi_\varepsilon(x) = f(x) + \varepsilon B(x).
\]
Note that, as $x$ tends to the boundary of $X$, $\psi_\varepsilon(x) \to +\infty$, so the barrier keeps the points inside $X$. As $\varepsilon \to 0$, $\psi_\varepsilon$ tends to $f$.

The function $B$ has the following properties:
\begin{itemize}
    \item $\operatorname{dom}(B) = \operatorname{int}(X)$;
    \item $B$ is convex;
    \item $B$ is smooth with
    \[
    \nabla B(x) = -\sum_{i=1}^{m} \frac{1}{g_i(x)}\, \nabla g_i(x),
    \]
    \[
    \nabla^2 B(x) = \sum_{i=1}^{m} \frac{1}{g_i(x)^2}\, \nabla g_i(x) \nabla g_i(x)^T + \sum_{i=1}^{m} \frac{1}{-g_i(x)}\, \nabla^2 g_i(x).
    \]
\end{itemize}

We now want to measure how far the solution of \eqref{eq:barrier} is from an optimal solution of (P). If $x_\varepsilon^*$ is an optimal solution of \eqref{eq:barrier}, then
\[
\nabla f(x_\varepsilon^*) + \sum_{i=1}^{m} \frac{\varepsilon}{-g_i(x_\varepsilon^*)}\, \nabla g_i(x_\varepsilon^*) = 0.
\]
Define
\[
\lambda_\varepsilon^* = \left( \frac{\varepsilon}{-g_1(x_\varepsilon^*)}, \dots, \frac{\varepsilon}{-g_m(x_\varepsilon^*)} \right) > 0.
\]
Consider the Lagrangian function associated with the given problem (P), $L(x, \lambda) := f(x) + \sum_{i=1}^{m} \lambda_i g_i(x)$. Then
\[
L(x, \lambda_\varepsilon^*) = f(x) + \sum_{i=1}^{m} (\lambda_\varepsilon^*)_i\, g_i(x)
\]
is convex and $\nabla_x L(x_\varepsilon^*, \lambda_\varepsilon^*) = 0$. Recall that (P) is a convex problem and strong duality holds, hence
\[
v(P) = \max_{\lambda \in \R^m_+} \inf_{x \in \R^n} L(x, \lambda).
\]
Consequently,
\[
v(P) \ge \min_x L(x, \lambda_\varepsilon^*) = L(x_\varepsilon^*, \lambda_\varepsilon^*).
\]
Finally, since $x_\varepsilon^*$ is feasible,
\[
f(x_\varepsilon^*) \ge v(P) \ge L(x_\varepsilon^*, \lambda_\varepsilon^*) = f(x_\varepsilon^*) + \sum_{i=1}^{m} (\lambda_\varepsilon^*)_i\, g_i(x_\varepsilon^*) = f(x_\varepsilon^*) - m\varepsilon,
\]
where $m\varepsilon$ is the \g{optimality gap}.

\begin{remark}
As $\varepsilon \to 0$, $f(x_\varepsilon^*) \to v(P)$.
\end{remark}
